\documentclass{ar2rc}
\usepackage{graphicx}
\usepackage[T1]{fontenc} % if needed
\usepackage{subfigure}
%\usepackage[dvipsnames]{color}
\definecolor{mygreen}{rgb}{.1,.75,.1}
\newcommand{\del}[1]{\textcolor{red}{\st{#1}}}
\newcommand{\comment}[1]{\textcolor{red}{#1}}
\newcommand{\add}[1]{\textcolor{mygreen}{#1}}
\title{Probing compressed Higgsinos at the FASER experiment}
\author{Shufang Su, Wei Su, Jin Min Yang, Pengxuan Zhu, Rui Zhu}
\doi{Manuscript ID: LN19434}
\begin{document}
\maketitle

\par The authors thank the reviewer for their critical assessment of our work. In the following, we address their concerns point by point. 
%\section*{Reviewer}

\RC The manuscript studies a compressed Higgsino scenario in the MSSM, focusing on the case in which the two Higgsino-like neutralinos become nearly degenerate and the next-to-lightest neutralino, $\tilde{\chi}_2^0$, becomes long-lived. The proposed signal is
\begin{equation*}
	\tilde{\chi}_2^0 \to \tilde{\chi}_1^0 \gamma, 
\end{equation*}
with detection at FASER/FASER-2. The paper argues that FASER-2 can probe a region with neutral Higgsino mass up to roughly 125–130 GeV and mass splitting
\begin{equation*}
	\Delta m \sim 4 - 30 {\rm MeV}, 
\end{equation*}
which is claimed to be complementary to conventional LHC Higgsino searches and chargino LLP searches.

\RCpar The topic is interesting, and the idea of probing long-lived neutral Higgsinos at FASER/FASER2 is potentially worthwhile. However, in my opinion the manuscript does not yet meet the standard expected for publication in a top-level journal. My main concern is that the central
claim relies on a quantitatively delicate prediction of an extremely small neutral Higgsino mass splitting, while the analysis does not treat this ingredient with the level of precision required by the claimed signal region.

\subsection*{Question 1}
\RC The key weakness of the manuscript is the treatment of the neutral Higgsino mass difference. The phenomenologically interesting region identified by the authors lies in the range $\Delta m_0 \sim 4-30$ MeV, but the analysis is based primarily on approximate tree-level expressions together with the qualitative statement that one-loop corrections mainly shift the cancellation condition without altering the overall conclusion. This is not sufficient for the present problem. The signal region is not simply a broadly compressed electroweakino spectrum; rather, it is a narrow cancellation-driven region in which the observable quantity is the small residual difference left after cancellation among much larger contributions. In such a situation, even modest loop corrections can shift $\Delta m_0$ by an amount comparable to, or larger than, the quoted signal window. Since the decay width is highly sensitive to the mass difference, the lifetime and therefore the predicted FASER event yield may change substantially even under a small shift in the splitting.

\RCpar For this reason, I do not find the manuscript’s treatment of loop effects adequate. In my view, the issue is not merely that the authors have not presented a full precision calculation; rather, the more basic problem is that, in such an extreme near-degenerate regime, even a nominal one-loop treatment may well be insufficient unless accompanied by a convincing assessment of missing higher-order effects and of the numerical stability of the spectrum itself. The manuscript assumes, without demonstrating it, that loop corrections merely shift the cancellation locus without materially affecting the phenomenologically relevant 4–30 MeV window. In such a cancellation-driven regime this assumption is not justified. Moreover, the accuracy of standard spectrum solvers in resolving splittings at the level of only a few MeV is itself open to question. A quantitatively reliable analysis would therefore require a substantially more careful treatment of the electroweakino spectrum, including at least a clearly defined loop-level calculation, a discussion of the dependence on renormalization prescription, a demonstration of numerical stability in the cancellation region, and a credible estimate of the uncertainty from uncomputed higher-order effects. Without such a treatment, the central reach claim does not appear sufficiently robust.

\AR Reply ONE. 

\subsection*{Question 2}
\RC A related concern is that the phenomenologically relevant region appears to arise only for a rather special choice of MSSM parameters, such that a strong cancellation occurs in the neutral Higgsino splitting. This means that the signal region corresponds to a narrow and highly tuned slice of parameter space rather than to a generic compressed Higgsino scenario. A tuned region is not automatically uninteresting, but if the authors wish to make a strong case for the broader importance of this scenario, they should either quantify the degree of tuning explicitly or provide a concrete motivation from a high-scale or SUSY-breaking perspective for why the relevant gaugino-mass relation should arise. At present, the manuscript does not do so, which weakens the broader phenomenological significance of the proposal.

\AR Reply TWO

\subsection*{Question 3}
\RC Independently of the issue above, I find the detector-level treatment of the proposed signal insufficient. The visible signature is the photon from $\tilde{\chi}_2^0 \to \tilde{\chi}_1^0 \gamma$. In the rest frame of $\tilde{\chi}_2^0$, the photon energy is approximately set by $\Delta m_0$, which in the claimed signal region is only 4–30 MeV. The actual observability of such a signal at FASER/FASER-2 depends on several nontrivial detector-level issues, including the laboratory-frame photon energy spectrum, the angular distribution of the photon, the probability of photon conversion, triggerability, photon reconstruction efficiency, and possible geometric losses. The manuscript demonstrates acceptance of the long-lived neutralino into the detector volume, but not the acceptance and reconstruction of the visible decay product that is supposed to establish the signal. This is a significant omission.

\AR 

\subsection*{Question 4}
\RC A further weakness is the absence of a serious treatment of backgrounds. The projected sensitivity appears to be based essentially on geometric acceptance, decay probability inside the fiducial volume, and an assumed three-event reach. This may be acceptable as a first exploratory estimate, but it is not sufficient for a strong phenomenological claim at the level targeted by the manuscript. For a mono-photon displaced signature, a realistic analysis should address at least the main potential background classes, such as neutrino-induced interactions, muoninduced secondaries or bremsstrahlung, and photon reconstruction or trigger limitations in the relevant soft-photon regime. In addition, the projected reach should be based on a detector-level signal definition with a quantified background model and a clear statistical procedure, rather than on an implicit three-event criterion. As it stands, the manuscript does not establish that the proposed mono-photon signature is experimentally realisable at FASER-2 with the claimed sensitivity, because neither the dominant backgrounds nor the photon efficient in the relevant energy regime are quantitatively demonstrated.

\AR \\

\subsection*{Summary}
\RC In summary, the manuscript presents an interesting idea: the use of FASER/FASER-2 to probe a very compressed neutral Higgsino sector in the MSSM through the radiative decay of a long-lived neutralino. However, the paper’s main conclusion rests on an extremely delicate prediction of the neutralino mass difference, and that prediction is not treated with the level of precision required by the claimed signal window. In particular, the mass-splitting calculation is not sufficiently robust for a signal region at the level of only a few MeV to a few tens of MeV, and the manuscript does not provide a realistic detector-level treatment of the photon signal or a quantitative estimate of the relevant backgrounds. For these reasons, I do not find the present analysis adequate for making a strong claim. \\

\AR 

\bibliography{ref_type1.bib}
\bibliographystyle{CitationStyle}

\end{document}
